% Generated draft: requires mathematical and source review.
\documentclass[11pt]{article}
\usepackage[margin=25mm]{geometry}
\usepackage{fontspec}
\setmainfont{Latin Modern Roman}
\usepackage{amsmath,amssymb,mathtools,unicode-math}
\setmathfont{Latin Modern Math}
\usepackage{enumitem,needspace,xurl,hyperref}
\setlength{\parindent}{0pt}
\setlength{\parskip}{5pt}
\setlength{\emergencystretch}{3em}
\newcommand{\N}{\mathbb{N}}
\newcommand{\Z}{\mathbb{Z}}
\newcommand{\Q}{\mathbb{Q}}
\newcommand{\R}{\mathbb{R}}
\newcommand{\C}{\mathbb{C}}
\newcommand{\F}{\mathbb{F}}
\newcommand{\A}{\mathbb{A}}
\newcommand{\PP}{\mathbb{P}}
\newcommand{\E}{\mathbb{E}}
\newcommand{\eps}{\varepsilon}
\newcommand{\dd}{\,\mathrm{d}}
\begin{document}
\section*{UnsolvedMath ID: 30003118}
\textbf{Source problem code:} catalogue entry
\subsection{Definitions and mathematical statement}
Fix an algebraic closure \(\overline{\mathbb Q}\subset\mathbb C\). For an integer \(d>2\), define
\[
\mathbb Q^{(d)}=\mathop{\mathrm{compositum}}_{\substack{\mathbb Q\subseteq K\subseteq\overline{\mathbb Q}\\ [K:\mathbb Q]\le d}}K,
\]
the smallest subfield of \(\overline{\mathbb Q}\) containing every such number field \(K\). Define its ring of integers by
\[
\mathcal O_{\mathbb Q^{(d)}}=\{\alpha\in\mathbb Q^{(d)}:\alpha\text{ satisfies a monic polynomial in }\mathbb Z[x]\}.
\]
For \(\alpha\in\overline{\mathbb Q}\), write its primitive irreducible polynomial over \(\mathbb Z\), with positive leading coefficient, as
\[
f_\alpha(x)=a\prod_{j=1}^{n}(x-\alpha_j),\qquad a>0.
\]
Its absolute logarithmic Weil height is
\[
h(\alpha)=\frac{1}{n}\left(\log a+\sum_{j=1}^{n}\log\max\{1,|\alpha_j|\}\right).
\]
A subset \(S\subseteq\overline{\mathbb Q}\) has Property \((N)\), also called the Northcott property, if for every real \(B\ge0\), the set \(\{\alpha\in S:h(\alpha)\le B\}\) is finite. Using the multiplicative height \(H(\alpha)=\exp(h(\alpha))\) gives the same property.

\textbf{Question.} Is it true that for every integer \(d>2\) and every real \(B\ge0\),
\[
\#\{\alpha\in\mathcal O_{\mathbb Q^{(d)}}:h(\alpha)\le B\}<\infty?
\]
Here the integer domain for \(d\) is the conventional interpretation of the supplied degree cutoff; the record writes only \(d>2\).

\textbf{Unresolved source qualification.} The displayed question follows the supplied original and cleaned statements. The supplied literature assessment instead uses the notation \(\mathcal O_{\mathbb Q_{\mathrm{tr}}^{(d)}}\), without defining \(\mathbb Q_{\mathrm{tr}}^{(d)}\) or identifying it with \(\mathbb Q^{(d)}\). Consequently, the intended relationship between that assessment and the displayed question remains unresolved. No totally real restriction has been added to the displayed formulation.
\subsection{Short English statement}
For each integer \(d>2\), form the compositum of all number fields of degree at most \(d\). The question asks whether only finitely many algebraic integers in this field have height below any prescribed bound. The supplied literature assessment uses a different, undefined field notation, so its applicability to this formulation requires review.
\subsection{Sources}
\begin{itemize}
\item Local UnsolvedMath problem registry, supplied complete record with permanent ID \(30003118\), source problem code \texttt{catalogue entry}, titled \emph{Northcott Property for Rings of Integers in Bounded-Degree Composita}; collection: Oberwolfach Reports Open Problems. No individual proposer is supplied; the proposed year is recorded as 2016.
\item Supplied report citation: \emph{Diophantische Approximationen (2017)}, Oberwolfach Reports, DOI 10.4171/owr/2016/21, pp.\ 1099--1187. Supplied source link: \url{https://doi.org/10.4171/owr/2016/21}. The record describes Oberwolfach Reports as documenting mathematical workshops and their open-problem discussions. The citation's 2017 label and its 2016 DOI component are retained as supplied.
\item Supplied literature source labeled ``OWR report'': \url{https://ems.press/content/serial-article-files/46625?nt=1}.
\item Supplied literature source labeled ``Bombieri--Zannier'': \url{https://arxiv.org/abs/1205.2491}. No further bibliographic details or specific supporting result are supplied.
\item Literature assessment dated exactly \texttt{2026-08-21}: the supplied curation labels the problem open and reports that the question for \(\mathcal O_{\mathbb Q_{\mathrm{tr}}^{(d)}}\) is explicitly open in the record and that no general affirmative result was located. This is supplied dated literature triage, not independently verified here; its time-sensitive status, cited claims, and applicability to \(\mathbb Q^{(d)}\) require independent verification.
\item The registry also supplies the statement label \texttt{corrected\_verified}, explained as preserving mathematical content while repairing notation, typography, or omitted definitions for self-contained publication. This is supplied metadata, not independent mathematical verification of this draft or its sources.
\end{itemize}
\end{document}
